% !TeX spellcheck = en_US
\documentclass[11pt]{article}
\usepackage{exerciseHead}

%%% ENCODING, FONT & LOCALIZATION
\usepackage[utf8]{inputenc} % use unicode input-encoding (allows unicode characters in the .tex files, default: Asci)
\usepackage{lmodern} % vectorized version of the "computer modern" font (default) (https://tex.stackexchange.com/q/1390/148164, https://tex.stackexchange.com/q/147194/148164)
\usepackage[ngerman, english]{babel} %% last language is main
%\shorthandoff{"} %babel breaks tikz-cd otherwise (activate if used)
\usepackage{csquotes} %% biblatex wants csquotes (localized quotes) if babel is present

%%% GENERAL IMPROVEMENTS
\usepackage[unicode, colorlinks]{hyperref} %% clickable links
\usepackage{enumitem} %% better (customizable) enumerations (e.g. \begin{enumerate}[label=(\roman*)])


%%% STYLING
% \setlength{\parindent}{0em} % No indentation for new paragraphs
\usepackage{emptypage} % removes page numbers on empty pages
% \makeatletter \@twosidefalse \makeatother
\usepackage{xcolor}
\usepackage[many]{tcolorbox}
\usepackage{fancyhdr}
\usepackage{mdframed}
\usepackage{booktabs}
\usepackage{multirow}

\usepackage{marvosym} % special letters


%%% GRAPHICS & PLOTS

% \usepackage{graphicx} % enable pictures
\usepackage{wrapfig} % text wrapping around figures
\usepackage{wrapstuff}
\usepackage{tikz, tikz-cd, pgfplots} % plots and sketches
\usetikzlibrary{bending,positioning,calc}
\usetikzlibrary{decorations.pathreplacing}
\usetikzlibrary{arrows.meta}
\pgfplotsset{compat=1.18}


\usepackage{subcaption}

%%% CODE
\usepackage[ruled,linesnumbered]{algorithm2e}

% fix algomathdisplay (see: https://tex.stackexchange.com/a/611426/148164)
\makeatletter
\renewenvironment{algomathdisplay}
 {\[}
 {\@endalgocfline\vspace{-\baselineskip}\]{\DontPrintSemicolon\;}}
\makeatother
%%% MATH
\usepackage{mathtools} % fixes some things from amsmath and adds some new features
\usepackage{amsthm, amssymb} % amsthm -> proof env, amssymb -> additional symbols
% \usepackage{thmtools} % restateable theorem environments
\usepackage{aligned-overset} % https://tex.stackexchange.com/questions/257529/overset-and-align-environment-how-to-get-correct-alignment
%% usage: \overset{something}&{=}
\usepackage{xfrac}

%%% Special Symbols
% \usepackage{bbm} % blackboard numbers: usage \mathbbm{1} (amssymb \mathbb{1} looks bad)
% \usepackage{cancel} % strike through equation parts to indicate they cancel out
% %% contradiction symbol
% \usepackage{wasysym} % lightning symbol
% \newcommand*{\contradiction}[1][]{\quad\text{{\LARGE\lightning} #1}\ }

%%% BIBLIOGRAPHY
\usepackage[numbers]{natbib}

% Editing
\usepackage[draft]{fixme}
\fxsetup{theme=color}

\usepackage{xsim}


%% AMS Theorem Environments
\theoremstyle{plain}% default
\newtheorem{proposition}{Proposition}[section]
\newtheorem{lemma}[proposition]{Lemma}
\newtheorem{corollary}[proposition]{Corollary}
\newtheorem{theorem}[proposition]{Theorem}
\newtheorem{fact}[proposition]{Fact}

\theoremstyle{definition}
\newtheorem{definition}[proposition]{Definition}
\newtheorem{example}[proposition]{Example}
\newtheorem{assumption}[proposition]{Assumption}

\theoremstyle{remark}
\newtheorem{remark}[proposition]{Remark}
%%

\newcommand*{\dims}{d}

\DeclarePairedDelimiter{\abs}{\lvert}{\rvert}
\DeclarePairedDelimiter{\Set}{\{}{\}}
\NewDocumentCommand{\Indicator}{s m}{%
  \IfBooleanTF{#1}
    {\mathbf{1}_{\left\{#2\right\}}} % starred: wrap in braces
    {\mathbf{1}_{#2}}                 % unstarred: simple subscript
}

% spaces
\newcommand*{\real}{\mathbb{R}}
\newcommand*{\nat}{\mathbb{N}}
\newcommand*{\integer}{\mathbb{Z}}
\newcommand*{\complex}{\mathbb{C}}
\newcommand*{\rational}{\mathbb{Q}}

% Topology
\newcommand*{\closure}[1]{\overline{#1}}

% Probability
\newcommand{\E}{\mathbb{E}}
\renewcommand{\Pr}{\mathbb{P}}
\DeclareMathOperator{\Var}{Var}
\DeclareMathOperator{\Cov}{Cov}
% \newcommand*{\normal}{\mathcal{N}}
\newcommand{\normal}[1]{{\color{blue}\boldsymbol{\mathcal{N}}}\left(#1\right)}
\newcommand*{\density}{\varphi}
\DeclareMathOperator{\iid}{iid}
\newcommand*{\given}{\mid}

\newcommand*{\Binomial}{\mathrm{Bin}}
\newcommand*{\Poisson}{\mathrm{Poi}}

% random functions
\newcommand*{\rf}{\mathbf{f}}

% Linear Algebra
\DeclareMathOperator{\sgn}{sgn}
\DeclareMathOperator{\Span}{span}
\newcommand{\identity}{\mathbb{I}}
\newcommand*{\transpose}{T}

% Optimization
\DeclareMathOperator*{\argmin}{arg\,min}
\DeclareMathOperator*{\argmax}{arg\,max}

% Machine Learning
\newcommand*{\cost}{J}
\newcommand*{\Cost}{\mathbf{J}}
\newcommand*{\loss}{\ell}
\newcommand*{\noise}{\varsigma}
\newcommand*{\param}{w}
\newcommand*{\model}{\phi}
\newcommand*{\X}{X}
\newcommand*{\Y}{Y}


% Numerical Analysis
\newcommand*{\bigO}{\mathcal{O}}


%% =========== Color (use RGB for digital, CMYK for print) ==========

% ----- Uni Mannheim Corporate Design -----
\definecolor{primary}{RGB}{0,48,86}
% \definecolor{primary}{cmyk}{1,0.6,0.1,0.6}

\colorlet{primary10}{primary!10}
\colorlet{primary25}{primary!25}
\colorlet{primary40}{primary!40}
\colorlet{primary55}{primary!55}
\colorlet{primary70}{primary!70}
\colorlet{primary85}{primary!85}

\definecolor{accent}{RGB}{179,182,185}
% \definecolor{accent}{cmyk}{35,25,25,0}

% ------------------------------------------

% --- convenient color access ---
\newcommand*{\colorPrimary}[2][]{{\color{primary#1} #2}}
\newcommand*{\colorAccent}[1]{{\color{accent} #1}}

 % AFTER the solution true/false macro !!

% \editmodetrue

\ifeditmode
  \NewDocumentEnvironment{solution}{O{Solution} +b}{\begin{proof}[#1] #2\end{proof}}{}
  \newtheorem{exercise}{Exercise}[section]
\else
  \usepackage{xsim}
  \DeclareExerciseEnvironmentTemplate{theorem}{%
    \par\addvspace{10pt}
    \noindent\textbf{%
      \XSIMmixedcase{\GetExerciseName}~\GetExerciseProperty{counter}%
    }{\normalfont%
      \GetExercisePropertyT{subtitle}{ (\PropertyValue)}}\textbf{. }
    \IfInsideSolutionF{%
      \GetExercisePropertyT{points}{%
        \marginpar{%
          \printgoal{\PropertyValue}%
          \GetExercisePropertyT{bonus-points}{+\printgoal{\PropertyValue}}%
          \,\IfExerciseGoalSingularTF{points}
            {\XSIMtranslate{point}}
            {\XSIMtranslate{points}}%
        }%
      }%
    }%
  }{\par\addvspace{\baselineskip}}

  \renewcommand\printsolutions{%
    \def\currentchapter{}%
    \def\currentsection{}%
    \def\lastchapter{}%
    \def\lastsection{}%
    \chapter*{Solutions\addcontentsline{toc}{chapter}{Solutions}}%
    
    \ForEachUsedExerciseByType{%
      \let\lastchapter\currentchapter
      \let\lastsection\currentsection
      \edef\currentchapter{\ExercisePropertyGet{##1}{##2}{chapter-value}}%
      \edef\currentsection{\ExercisePropertyGet{##1}{##2}{section-value}}%
      % \ifx\lastchapter\currentchapter\else
      %   \addcontentsline{toc}{chapter}{Solutions Chapter \ExercisePropertyGet{##1}{##2}{chapter}}
      %   \chapter*{Solutions Chapter \ExercisePropertyGet{##1}{##2}{chapter}}
      % \fi
      \ifx\lastsection\currentsection\else
        \section*{Exercises \ExercisePropertyGet{##1}{##2}{section}}%
        \addcontentsline{toc}{section}{Exercises \ExercisePropertyGet{##1}{##2}{section}}%
      \fi
      \XSIMprint{solution}{##1}{##2}%
    }%
  }
  \xsimsetup{
      exercise/template=theorem,
      solution/template=theorem,
      exercise/the-counter = \thesection.\arabic{exercise},
      exercise/within=section,
      % solution/print=true,
  }
  % \SetExerciseParameters{exercise}{counter=theorem}
\fi
  \xsimsetup{
      exercise/the-counter = \arabic{exercise},
    %   solution/print=true,
  }


% Meta Information %%%%%%%%%%%%%%%%%%%%%%%%%%%%%%%%%%%

% included in all sheets (-> only one update per semester)
\course{Compl\'ements de probabilit\'es et statistiques}
\semester{Spring 2026}
\author{Felix Benning}

 % (file) for one update per semester
\sheetNumber{2}

%%%% Document %%%%%%%%%%%%%%%%%%%%%%%%%%%%%%%%%%%%%%%%
\begin{document}

\maketitle

\textbf{
The goal is to have something to talk about in the oral exam.  Skip an exercise
if you get stuck. That is why there are many.
% I am more interested in how you
% would solve them than the actual solution, i.e.\ do not prioritize tedious
% calculations.
}

% 
\begin{exercise}[subtitle={Properties of the cumulant generating function}]
    \label{ex: properties cumulant generating function}
    Let \(X\) be a random variable. Then the cumulant generating function
    \(\CGF_X(t) = \log \MGF_X(t)\) has the following properties:
    \begin{enumerate}[label=(\roman*), noitemsep]
        \item\label{it: cgf 0=0} \(\CGF_X(0) = 0\).
        \item\label{it: cgf linear} \(\CGF_{aX + b}(t) = \scp{b, t} + \CGF_X(\scp{a, t})\) for any \(a \in \real\), \(b,t\in \hilbert\).
        \item\label{it: cumulant independent sums} Let \(Y\) be \underline{independent} from \(X\), then
        \[
            \CGF_{X+Y}(t) = \CGF_X(t) + \CGF_Y(t)
        \]
    \end{enumerate}
    If the moment generating function is finite in a neighborhood of zero,
    then \(\CGF_X\) is infinitely often differentiable and
    \begin{enumerate}[label=(\roman*), resume, noitemsep]
        \item \(\CGF_X'(0) = \E{X}\),
        \item \(\CGF_X''(0) = \var{X}\)
    \end{enumerate}
\end{exercise}

\begin{solution}
    For \ref{it: cgf 0=0}-\ref{it: cumulant independent sums} follow from the
    corresponding properties of the MGF (Exercise \ref{ex: properties moment
    generating functions}) by application of the logarithm.  For example,
    \(\CGF_X(0)=0\) follows directly from \(\MGF_X(0)=1\). For convexity
    let \(t, s\in \hilbert\) and \(\lambda \in [0,1]\). 

    The derivatives are straightforward calculations using the derivatives of the MGF
    and the chain rule due to \(\CGF_X(t) = \log \MGF_X(t)\). That is
    \begin{alignat*}{3}
        \CGF_X'(t)
        &= \frac{\MGF_X'(t)}{\MGF_X(t)}
        \qquad &\implies &\CGF_X'(0) = \E{X},
        \\
        \CGF_X''(t)
        &= \frac{\MGF_X''(t)}{\MGF_X(t)} - \frac{(\MGF_X'(t))^2}{(\MGF_X(t))^2}
        \qquad &\implies &\CGF_X''(0)
        \begin{aligned}[t]
        &= \E{X^2} - (\E{X})^2
        \\
        &= \var{X}.
        \end{aligned}
    \end{alignat*}
    This proves the claims.
\end{solution}



\begin{exercise}[subtitle={Almost sure but not \(L^p\) convergence}]
    \label{ex: as but not Lp}
    Let \(U\sim \uniform(0,1)\) and define
    \[
        X_n = 2^n \ind{[0, 2^{-n})}(U).
    \]
    Show that \(X_n \overset{\as}\to 0\) but \(X_n\) does not converge to \(0\)
    in \(L^1\) and hence in no other \(L^p\) space with \(p \ge 1\).
\end{exercise}

\begin{solution}
    First, we show that \(X_n \overset{\as}\to 0\). For any \(\epsilon > 0\)
    we have
    \[\begin{aligned}
        \sum_{n=0}^\infty \prob{\abs{X_n - 0} \ge \epsilon}
        &\le \sum_{n=0}^\infty \prob{X_n \ge \min\set{\epsilon, 1}}
        \\
        &= \sum_{n=0}^\infty\prob{U < 2^{-n}} = \sum_{n=0}^\infty 2^{-n} < \infty.
    \end{aligned}
    \]
    Consequently we have \(X_n \overset{\as}\to 0\) by the characterization of almost sure convergence.
    To see that \(X_n\) does not converge to \(0\) in \(L^1\) we observe
    \[
        \E{\abs{X_n - 0}}
        = \E{X_n}
        = 2^n \prob{U < 2^{-n}} = 1
        \not\to 0.
        \qedhere
    \]
\end{solution}


\begin{exercise}[subtitle={Products of random variables}]
    For \((X_n)_{n\in \nat}\) independent and identically distributed random
    variables with \(X_n > 0\) and \(\E{|\ln(X_1)|} < \infty\) prove that the
    `geometric mean' converges almost surely to the `geometric expectation',
    that is'
    \[
        \Bigl(\prod_{i=1}^n X_i \Bigr)^\frac1n
        \overset{\as}\to \exp\bigl(\E[\big]{\ln(X_1)}\bigr).
    \]
\end{exercise}

\begin{solution}
    Observe that the logarithm of the geometric mean converges by the law of
    large numbers to the expectation of the logarithm, that is
    \[
        \ln\Bigl(\Bigl(\prod_{i=1}^n X_i \Bigr)^\frac1n\Bigr)
        = \frac1n \sum_{i=1}^n \ln(X_i)
        \overset{\as}\to \E[\big]{\ln(X_1)}.
    \]
    Since \(\exp(x)\) is continuous the claim follows from continuous mapping.
\end{solution}

\begin{exercise}[subtitle={Properties of cumulants}]
    \label{ex: properties cumulants}
    Prove that cumulants have the following properties: 
    \begin{enumerate}[label=(\alph*), noitemsep]
        \item\label{it: (cumulants) accumulate} \textbf{Accumulation.}
        If \(X\) and \(Y\) are independent, then
        the cumulants ``accumulate''
        \[
            \cumulant_n[X+Y] = \cumulant_n[X] + \cumulant_n[Y].
        \]

        \item\label{it: (cumulants) homogeneity}\textbf{Homogeneity.} For any \(a \in \real\) we have
        \[
            \cumulant_n[aX] = a^n \cumulant_n[X].
        \]

        \item\label{it: (cumulants) constant} Prove that the cumulant generating function of a constant random variable \(Y=\mu\)
        is given by
        \[
            \CGF_Y(t) = \mu t.
        \]
        For a random variable \(X\) with \(\E{X} = \mu\) deduce that
        \(\cumulant_k[X-\mu] = \cumulant_k[X]\) for all \(k \ge 2\) and
        \(\cumulant_1[X-\mu]=0\).
    \end{enumerate}
\end{exercise}

\begin{solution}
\begin{enumerate}[wide=0pt]
    \item[\ref{it: (cumulants) accumulate}]
    We have by the sum property for the cumulant generating function
    \begin{align*}
        \sum_{n=1}^\infty \cumulant_n[X+Y] \frac{t^n}{n!}
        &= \CGF_{X+Y}(t)
        \\
        &= \CGF_X(t) + \CGF_Y(t)
        = \sum_{n=1}^\infty \cumulant_n[X] \frac{t^n}{n!} + \sum_{n=1}^\infty \cumulant_n[Y] \frac{t^n}{n!}.
        \\
        &= \sum_{n=1}^\infty \bigl(\cumulant_n[X] + \cumulant_n[Y]\bigr) \frac{t^n}{n!}.
    \end{align*}
    A comparison of the coefficients proves the claim.

    \item[\ref{it: (cumulants) homogeneity}]
    By the scaling property of the cumulant generating function we have
    \[
        \sum_{n=1}^\infty \cumulant_n[aX] \frac{t^n}{n!}
        = \CGF_{aX}(t)
        = \CGF_X(at)
        = \sum_{n=1}^\infty \cumulant_n[X] \frac{(at)^n}{n!}
        = \sum_{n=1}^\infty a^n \cumulant_n[X] \frac{t^n}{n!}
    \]
    and the claim follows again by comparing the coefficients.

    \item[\ref{it: (cumulants) constant}] We have
    for a constant random variable \(Y=\mu\) that
    \[
        \MGF_Y(t) = e^{t \mu} \implies \CGF_Y(t) = t \mu.
    \]
    Comparing coefficients with the series expansion we consequently have
    \(\cumulant_1[Y] = \mu\) and \(\cumulant_k[Y] = 0\) for \(k \ge 2\).
    Since a constant random variable is independent of every other random variable,
    we have by the accumulation property that
    \[
        \cumulant_k[X-\mu]
        = \cumulant_k[X] + \cumulant_k[-\mu]
        = \cumulant_k[X]
        \quad\text{for } k \ge 2,
    \]
    and \(\cumulant_1[X-\mu] = \cumulant_1[X] + \cumulant_1[-\mu] = 0\).
    \qedhere
    \end{enumerate}
\end{solution}

\begin{exercise}[subtitle={Cumulant CLT}]
    \label{ex: cumulant CLT}
    Let \(X_1, X_2, \ldots\) be independent copies of a random variable
    \(X\) with mean \(\mu\), variance \(\sigma^2\) and finite moments of all orders.
    Let \(\mean{X}_n = \frac1n \sum_{i=1}^n X_i\) be the empirical mean.
    Show that the cumulants \(\cumulant_k[\sqrt{n}(\mean{X}_n - \mu)]\)
    of the random variable \(\sqrt{n}(\mean{X}_n - \mu)\) satisfy
    \begin{align*}
        \cumulant_1[\sqrt{n}(\mean{X}_n - \mu)] &= 0,
        \\
        \cumulant_2[\sqrt{n}(\mean{X}_n - \mu)] &= \sigma^2,
        \\
        \cumulant_k[\sqrt{n}(\mean{X}_n - \mu)] &\to 0
        \quad\text{for } k \ge 3.
    \end{align*}
    Conclude that the cumulant generating function of
    \(\sqrt{n}(\mean{X}_n - \mu)\) converges to the cumulant generating function
    of the normal distribution \(\normal{0, \sigma^2}\).
\end{exercise}

\begin{solution}
    The first and second cumulant are the expectation and variance
    respectively and therefore
    \begin{align*}
        \cumulant_1[\sqrt{n}(\mean{X}_n - \mu)]
        &= \E{\sqrt{n}(\mean{X}_n - \mu)} = 0,
        \\
        \cumulant_2[\sqrt{n}(\mean{X}_n - \mu)]
        &= \var{\sqrt{n}(\mean{X}_n - \mu)}
        = n \var{\mean{X}_n}
        = n \tfrac{\sigma^2}{n} = \sigma^2.
    \end{align*}
    For the higher order cumulants we have we use the accumulation property
    and homogeneity of Exercise \ref{ex: properties cumulants} to get
    \begin{align*}
        \cumulant_k[\sqrt{n}(\mean{X}_n - \mu)]
        &= \cumulant_k\Bigl[\frac1{\sqrt{n}} \sum_{i=1}^n (X_i - \mu)\Bigr]
        \\
        \overset{\text{accum.}}&= n\, \cumulant_k[\tfrac1{\sqrt{n}} (X_1 - \mu)]
        \\
        \overset{\text{hom.}}&= n^{1-\frac{k}{2}} \cumulant_k[X_1 - \mu].
        \\
        &= n^{1-\frac{k}{2}} \cumulant_k[X_1] \qquad \forall k \ge 2.
    \end{align*}
    Clearly for \(k\ge 3\) the cumulants decay to zero. Convergence of the
    cumulant generating function then follows by dominated convergence since the
    series of analytic functions converges absolutely in a neighborhood of zero.
\end{solution}


\begin{exercise}[subtitle={Monte Carlo}]
    \label{ex: monte carlo}
    In this exercise we will devise a way to compute the integral
    \[
        I \coloneq \int_{[0,1]^\dims} f(x) dx
    \]
    of a continuous function \(f\colon [0,1]^\dims \to \real\).  Let
    \((U_n)_{n\in \nat}\) be independent and uniformly distributed on
    \([0,1]^\dims\).
    \begin{enumerate}[label=(\alph*), noitemsep]
        \item\label{it: (monte carlo) convergence}
        Prove that \(\hat I_n \coloneq \frac1n\sum_{i=1}^n f(U_i) \overset\as\to I\).

        \item\label{it: (monte carlo) PAC bound}
        Using \(\norm{f}_\infty = \sup_{x\in [0,1]^\dims} \abs{f(x)}\) Prove that
        \[
            \prob[\big]{\abs{\hat I_n - I} > \epsilon} \le 2\exp\Bigl(-\frac{n \epsilon^2}{2\norm{f}_\infty^2}\Bigr)
        \]

        \item\label{it: (monte carlo) sample complexity}
        Let \(n = n(\epsilon, \delta)\) be the smallest \(n\) such that the following
        PAC bound holds
        \[
            \prob[\big]{\abs{\hat I_n - I} > \epsilon} \le \delta
        \]
        Prove that
        \[
            n(\epsilon, \delta)
            \le \ceil[\Big]{2\norm{f}_\infty^2\frac{\ln(\tfrac2\delta)}{\epsilon^2}}
            \in \bigO\Bigl(\norm{f}_\infty^2\frac{\log(\frac1\delta)}{\epsilon^2}\Bigr)
        \]
    \end{enumerate}
\end{exercise}

\begin{solution}
\begin{enumerate}[wide=0pt]
    \item[\ref{it: (monte carlo) convergence}]
    Since \(I= \E{f(U)}\) this is the strong law of large numbers.

    \item[\ref{it: (monte carlo) PAC bound}]
    Since \(f\) is continuous, it is bounded on the compact set \([0,1]^\dims\)
    and thereby \(f(U) \in [-\norm{f}_\infty, \norm{f}_\infty]\).
    By Hoeffding's inequality 
    for bounded random variables we therefore have the claim.

    \item[\ref{it: (monte carlo) sample complexity}]
    By the previous exercise it is sufficient to have
    \begin{align*}
        2\exp\Bigl(-\frac{n\epsilon^2}{2\norm{f}_\infty^2}\Bigr) \le \delta
        &\iff -\frac{n\epsilon^2}{2\norm{f}_\infty^2} \le \ln(\tfrac{\delta}2)
        \\
        &\iff \ln(\tfrac{2}\delta) \le \frac{n \epsilon^2}{2\norm{f}_\infty^2}
        \\
        &\iff 2\norm{f}_\infty^2 \frac{\ln(\tfrac2\delta)}{\epsilon^2} \le n
    \end{align*}
    Since this is a sufficient condition, the smallest \(n(\delta, \epsilon)\)
    is smaller than the first natural number that satisfies it, leading
    to the claim.
    \qedhere
\end{enumerate}
\end{solution}


\end{document} 

%%% Local Variables: 
%%% mode: latex
%%% TeX-master: t
%%% End: 

